\documentclass[british,aps,prl,superscriptaddress,nofootinbib,times,reprint]{revtex4-1}
\usepackage[T1]{fontenc}
\usepackage[latin9]{inputenc}
\usepackage{geometry}
\geometry{verbose,tmargin=2cm,bmargin=2cm,lmargin=2cm,rmargin=2cm,headheight=2cm,headsep=2cm,footskip=2cm}
\setcounter{secnumdepth}{3}
\usepackage{verbatim}
\usepackage{refstyle}
\usepackage{amsmath}
\usepackage{amsthm}
\usepackage{amssymb}
\usepackage{esint}
\makeatletter
%%%%%%%%%%%%%%%%%%%%%%%%%%%%%% LyX specific LaTeX commands.
\AtBeginDocument{\providecommand\thmref[1]{\ref{thm:#1}}}
\AtBeginDocument{\providecommand\lemref[1]{\ref{lem:#1}}}
\AtBeginDocument{\providecommand\propref[1]{\ref{prop:#1}}}
\AtBeginDocument{\providecommand\eqref[1]{\ref{eq:#1}}}
\AtBeginDocument{\providecommand\figref[1]{\ref{fig:#1}}}
\AtBeginDocument{\providecommand\secref[1]{\ref{sec:#1}}}
\AtBeginDocument{\providecommand\ssecref[1]{\ref{ssec:#1}}}
\AtBeginDocument{\providecommand\defnref[1]{\ref{defn:#1}}}
\AtBeginDocument{\providecommand\tblref[1]{\ref{tbl:#1}}}

\newref{tbl}{name=Table~}{}
\newref{defn}{name=Definition~}{}
\newref{sec}{name=Section~}
\newref{ssec}{name=Subsection~}

\RS@ifundefined{subref}
  {\def\RSsubtxt{section~}\newref{sub}{name = \RSsubtxt}}
  {}
\RS@ifundefined{thmref}
  {\def\RSthmtxt{Theorem~}\newref{thm}{name = \RSthmtxt}}
  {}
\RS@ifundefined{lemref}
  {\def\RSlemtxt{Lemma~}\newref{lem}{name = \RSlemtxt}}
  {}
\RS@ifundefined{propref}
  {\def\RSproptxt{Proposition~}\newref{prop}{name = \RSproptxt}}
  {}

%%%%%%%%%%%%%%%%%%%%%%%%%%%%%% Textclass specific LaTeX commands.
\theoremstyle{plain}
\newtheorem{thm}{\protect\theoremname}

\theoremstyle{plain}
\newtheorem*{thm*}{\protect\theoremname}

  \theoremstyle{definition}
  \newtheorem{defn}{\protect\definitionname}

  \theoremstyle{remark}
  \newtheorem{rem}{\protect\remarkname}

  \theoremstyle{remark}
  \newtheorem{defnrem}{\protect\remarkname}[defn]


  \theoremstyle{remark}
  \newtheorem{thmrem}{\protect\remarkname}[thm]

  \theoremstyle{plain}
  \newtheorem*{prop*}{\protect\propositionname}
  \theoremstyle{plain}
  \newtheorem{lem}[thm]{\protect\lemmaname}
  \theoremstyle{plain}
  \newtheorem{prop}[thm]{\protect\propositionname}
  \theoremstyle{definition}
  \newtheorem{example}[thm]{\protect\examplename}
  \theoremstyle{definition}
  \newtheorem*{notn*}{\protect\notationname}


%%%%%%%%%%%%%%%%%%%%%%%%%%%%%% User specified LaTeX commands.
\usepackage{paralist}
\usepackage{graphicx}
\graphicspath{ {images/} }

\usepackage{xcolor}
%\usepackage{url}
\usepackage{hyperref}
\hypersetup{
    colorlinks,
    linkcolor={red!50!black},
    citecolor={blue!50!black},
    urlcolor={blue!80!black}
}

\makeatother

\usepackage{babel}
  \providecommand{\definitionname}{Definition}
  \providecommand{\examplename}{Example}
  \providecommand{\lemmaname}{Lemma}
  \providecommand{\propositionname}{Proposition}
  \providecommand{\remarkname}{Remark}
  \providecommand{\theoremname}{Theorem}
  \providecommand{\notationname}{Notation}
\setlength{\itemindent}{0pt}

% %%%%%%%%%%%%%%%%%%%%%%%%%%%%%%%%%%This stuff is required for transposing matrices | fine I didn't want to do donkey work
% \usepackage{xparse}
% \usepackage{environ}

% \ExplSyntaxOn
% \NewEnviron{bmatrixT}
%  {
%   \marine_transpose:V \BODY
%  }

% \int_new:N \l_marine_transpose_row_int
% \int_new:N \l_marine_transpose_col_int
% \seq_new:N \l_marine_transpose_rows_seq
% \seq_new:N \l_marine_transpose_arow_seq
% \prop_new:N \l_marine_transpose_matrix_prop
% \tl_new:N \l_marine_transpose_last_tl
% \tl_new:N \l_marine_transpose_body_tl

% \cs_new_protected:Nn \marine_transpose:n
%  {
%   \seq_set_split:Nnn \l_marine_transpose_rows_seq { \\ } { #1 }
%   \int_zero:N \l_marine_transpose_row_int
%   \prop_clear:N \l_marine_transpose_matrix_prop
%   \seq_map_inline:Nn \l_marine_transpose_rows_seq
%    {
%     \int_incr:N \l_marine_transpose_row_int
%     \int_zero:N \l_marine_transpose_col_int
%     \seq_set_split:Nnn \l_marine_transpose_arow_seq { & } { ##1 }
%     \seq_map_inline:Nn \l_marine_transpose_arow_seq
%      {
%       \int_incr:N \l_marine_transpose_col_int
%       \prop_put:Nxn \l_marine_transpose_matrix_prop
%        {
%         \int_to_arabic:n { \l_marine_transpose_row_int }
%         ,
%         \int_to_arabic:n { \l_marine_transpose_col_int }
%        }
%        { ####1 }
%      }
%    }
%    \tl_clear:N \l_marine_transpose_body_tl
%    \int_step_inline:nnnn { 1 } { 1 } { \l_marine_transpose_col_int }
%     {
%      \int_step_inline:nnnn { 1 } { 1 } { \l_marine_transpose_row_int }
%       {
%        \tl_put_right:Nx \l_marine_transpose_body_tl
%         {
%          \prop_item:Nn \l_marine_transpose_matrix_prop { ####1,##1 }
%          \int_compare:nF { ####1 = \l_marine_transpose_row_int } { & }
%         }
%       }
%      \tl_put_right:Nn \l_marine_transpose_body_tl { \\ }
%     }
%    \begin{bmatrix}
%    \l_marine_transpose_body_tl
%    \end{bmatrix}
%  }
% \cs_generate_variant:Nn \marine_transpose:n { V }
% \cs_generate_variant:Nn \prop_put:Nnn { Nx }
% \ExplSyntaxOff


\begin{document}
\title{A non-contextual hidden variable model for quantum mechanics}
\author{Atul Singh Arora}
\email{ms11003@iisermohali.ac.in}
\selectlanguage{british}%
\affiliation{Indian Institute of Science Education and Research (IISER), Mohali}
\author{Arvind}
\email{arvind@iisermohali.ac.in}
\selectlanguage{british}%
\affiliation{Indian Institute of Science Education
and Research (IISER), Mohali}
\begin{abstract}
We propose a non-contextual hidden variable model,
consistent with all predictions of quantum
mechanics (QM).  A careful scrutiny of
consistency requirements between any hidden
variable model and  quantum mechanics and the
corresponding no-go theorems, leads us to the
conclusion that the notion of contextuality is not
a necessary feature of QM.  The alternative view
that emerges, hinges on identifying a new classical
property, which we call ``multiplicativity''.  It turns
out that, by relaxing this condition of
``multiplicativity'' on hidden variable models, they
can be made  consistent with QM.  Advantages of
this view are illustrated by considering the
implications of non-multiplicativity to locality,
entanglement, superposition, and  apparent
contextuality. 
\end{abstract}
\pacs{}
\maketitle

Quantum mechanics has been one of the most
successful theories in physics so far, however,
there has not yet been a final word on its
completeness and interpretation~\cite{??}.
Einstein's~\cite{EinsteinEPR} work on the
incompleteness of QM and the subsequent seminal
work of Bell~\cite{Bell1964}, assessing the
compatibility of a more complete model involving
hidden variables (HV) and locality with QM, has
provided deep insights into  how the quantum world
differs from its classical counterpart.  In recent
times, these insights have been of pragmatic
utility in the areas of quantum information
processing, where EPR pairs are fundamental motifs
of entanglement~\cite{Ekert,PironioRndmnssCrtfcn}.
The work of  Kochen Specker~\cite{KochenSpecker}
and Gleason et.
al.~\cite{Gleason,BellOnHiddenVariables,Peres,Mermin})
broadened the schism between  HV models and QM.
They showed that it was contextuality and not
non-locality which was at the heart of this schism
and the incompatibility between HV models and QM
can arise even for a single quantum system.
Contextuality has thus been identified as a
fundamental non-classical feature of the quantum
world and experiments have also been proposed and
conducted to this effect~\cite{SimonContExpProp,
HuangContExp,YangContExp,HasegawaContExp}.
Contextuality, on the one hand has led to
investigations on the foundational aspects of
QM~\cite{PawelCntxClsscl,?????}, and on the other
hand has been harnessed for computation and
cryptography~\cite{HowardCntxCmptn,CabelloCntxScrt,??}.
Quantitative measures of contextuality based on
memory~\cite{MatthiasCntxMmry} have also been used
to demonstrate the completeness of
QM~\cite{CabelloMmryQM}.  However, not all HV
models are incompatible with QM, and a very
important case in point is  Bohm's model based on
precise trajectories for quantum
systems~\cite{Bohm1,Bohm2}. 


A closer look at the aforesaid no-go theorems
reveals that contextuality is in fact not a
necessary feature of QM (see \figref{block}).
While this has been shown earlier by taking
disturbances due to measurements into
consideration~\cite{NoContextuality,LaCourNoCntx},
we take a different approach and construct a
HV model consistent with QM.  The predictions of
the aforesaid model match with those of QM for systems
described by arbitrary finite-dimensional Hilbert
spaces.  \begin{figure}[h]
\includegraphics[width=0.9\columnwidth]{block1}
\caption{Non-contextuality is not inconsistent
with QM} \label{fig:block}\end{figure} The
algebraic constraints obeyed by quantum
observables may not necessarily be obeyed by the
predictions of an HV model at the level of
individual outcomes.  Those HV models that obey
such algebraic constraints possess a special
property which we call `multiplicativity'.  
Enlarging the scope of HV models to include
non-multiplicative models, permits us to construct an
explicit non-contextual model consistent with QM.
The re-examination of the very proof of
`contextuality'  is used to illustrate
non-multiplicativity, an intrinsic quantum feature, 
which we propose as a possible alternative to contextuality. 


%{\noindent \bf Notation:} 
\begin{notn*}
(a) $\psi\in\mathcal{H}$ 
represents a pure quantum mechanical
state of the system in the Hilbert space
$\mathcal{H}$, (b) $\hat{\mathcal{H}}$ is defined
to mean the set of Hermitian 
observables for the system, (c)
$[\mathcal{H}]$ is defined to mean
$(\mathcal{H},\,\mathbb{R}^{\otimes})$, which
represents the state of the system including HVs,
(d) $[\psi]\in[\mathcal{H}]$ will represent the
state of the system, including HVs, (e) a
prediction map is
$M:\{ \hat{\mathcal{H}},[\mathcal{H}] \}\to\mathbb{R}$,
(f) a sequence map is
$S:\{ \hat{\mathcal{H}},[\mathcal{H}],\mathbb{R} \}\to[\mathcal{H}]$,
(g) $f$ is an arbitrary map from $\{
\hat{\mathcal{H}},\hat{\mathcal{H}},\dots\hat{\mathcal{H}}
\} \to \hat{\mathcal{H}}$ constructed using
multiplication and addition of the observables,
and multiplication with complex
numbers\footnote{Strictly, the map $f$ would
depend on the operators to ensure preservation of
Hermiticity; thus a given $f$ will be defined only
for a subset of $\{
\hat{\mathcal{H}},\hat{\mathcal{H}},\dots\hat{\mathcal{H}}
\}$.}, (h) $\tilde{f}$ is a map constructed by
replacing observables in $f$ with real numbers.
\end{notn*}

\begin{defn} A theory is non-contextual, if it
provides a map $M: \{
\hat{\mathcal{H}},[\mathcal{H}] \} \to\mathbb{R}$
to explain measurement outcomes. A theory which is
not non-contextual is contextual.\end{defn}


\begin{defnrem} The notion of non-contextuality may be
extended to maps. A map is non-contextual if it is
of the form $M: \{ \hat{\mathcal{H}},[\mathcal{H}]
\} \to\mathbb{R}$; prediction maps as defined above, for individual observables, are non-contextual by construction.  
\end{defnrem}
\begin{defnrem}
Broader definitions in the literature have been
suggested which declare a larger set of theories
as non-contextual. For our purposes however, this
restricted definition will suffice. The term
contextual is used to suggest that the value an
operator takes might depend on which other
compatible observable it is being  measured with.
\end{defnrem}


Before we re-examine the result of Kochen and
Specker~\cite{KochenSpecker} that non-contextual
models cannot be compatible with quantum
mechanics and show that multiplicativity can provide an
alternative to contextuality, we define
multiplicativity more precisely.
\begin{defn}
A prediction map $M$ is \emph{multiplicative} iff
\begin{align*} &
M(f(\hat{B}_{1},\hat{B}_{2},\dots\hat{B}_{N}),[\psi])
= \\ & \tilde
f(M(\hat{B}_{1},[\psi]),M(\hat{B}_{2},[\psi]),\dots
M(\hat{B}_{N},[\psi])), \end{align*}
where $\hat{B}_{i}\in\hat{\mathcal{H}}$ are
arbitrary mutually commuting observables and
$[\psi]\in[\mathcal{H}]$. A
\emph{non-multiplicative} map is one that is not
multiplicative.  
\end{defn} 
Note that if $M$ is taken to represent the
measurement outcome (in QM), then for states of
the system which are simultaneous eigenkets of
$\hat{B}_{i}$s, $M$ must clearly be
multiplicative. It is, however, not obvious that
this property must always hold. For example,
consider two spin-half particles in the state 
$\left|1\right\rangle \otimes\left|1\right\rangle
$ written in the computational basis and take
operators 
$\hat{B}_{1}=\hat{\sigma}_{x}\otimes\hat{\sigma}_{x}$,
$\hat{B}_{2}=\hat{\sigma}_{y}\otimes\hat{\sigma}_{y}$
and
$\hat{C}=\hat{B}_{1}\hat{B}_{2}=-\hat{\sigma}_{z}\otimes\hat{\sigma}_{z}$
written in terms of Pauli operators.
We must have $M(\hat{C})=-1$ while
$M(\hat{B}_{1})=\pm1$ and $M(\hat{B}_{2})=\pm1$
independently, according to QM, with probability
half. Here multiplicativity does not seem to hold.
Antithetically, it is clear that if one first
measures $\hat{B}_{1}$ and subsequently measures
$\hat{B}_{2}$, then the product of the results
must be $-1$. This is consistent with measuring
$\hat{C}$. To capture this property we define
\emph{sequential multiplicativity} as follows.
\begin{defn} 
A prediction map $m$ is
\emph{sequentially multiplicative} for a given
sequence map $s$, iff \begin{align*}
&M(f(\hat{B}_{1},\hat{B}_{2},\dots\hat{B}_{N}),[\psi_{1}])=\\
&\tilde
f(M(\hat{B}_{1},[\psi_{k_{1}}]),M(\hat{B}_{2},[\psi_{k_{2}}]),\dots,M(\hat{B}_{N},[\psi_{k_{N}}])),
\end{align*} 
where
$\mathbf{k}=(k_{1},k_{2},\dots,k_{N})\in\{
N! {\rm ~permutations~of~}  k'{\rm s}
\}$, $\hat{B}_{i}\in\hat{\mathcal{H}}$ are
arbitrary mutually commuting observables,
$[\psi_{i}]\in[\mathcal{H}]$ and
$[\psi_{k+1}] :=
S(\hat{B}_{k},[\psi_{k}],M(\hat{B}_{k},[\psi_{k}]))$,
$\forall\,[\psi_{i}]$.  
\label{defn:seqnMltpl}\end{defn} 
With these definitions we are now ready to discuss
the `proof of contextuality'. We first state the
contextuality theorem in our notation:
\begin{thm*} Let a map
$M:\hat{\mathcal{H}}\to\mathbb{R}$, be s.t. (a)
$M(\hat{\mathbb{I}})=1$, (b)
$M(f(\hat{B}_{1},\hat{B}_{2},\dots))=\tilde
f(M(\hat{B}_{1}),M(\hat{B}_{2}),\dots)$, for any
arbitrary function $f$, where $\hat{B}_{i}$ are
mutually commuting Hermitian operators. If $m$ is
assumed to describe the outcomes of measurements,
then no $M$ exists which is consistent with all
predictions of QM. 
\label{thm:KS}
\end{thm*}

\begin{proof} Peres Mermin [PM]
($\left|\mathcal{H}\right|\ge4$)
\cite{Peres,Mermin}:
For a system composed of two spin-half particles 
consider the following set of operators \[
\hat{A}_{ij}\doteq\left[\begin{array}{ccc}
\hat{\mathbb{I}}\otimes\hat{\sigma}_{x} &
\hat{\sigma}_{x}\otimes\hat{\mathbb{I}} &
\hat{\sigma}_{x}\otimes\hat{\sigma}_{x}\\
\hat{\sigma}_{y}\otimes\hat{\mathbb{I}} &
\hat{\mathbb{I}}\otimes\hat{\sigma}_{y} &
\hat{\sigma}_{y}\otimes\hat{\sigma}_{y}\\
\hat{\sigma}_{y}\otimes\hat{\sigma}_{x} &
\hat{\sigma}_{x}\otimes\hat{\sigma}_{y} &
\hat{\sigma}_{z}\otimes\hat{\sigma}_{z}
\end{array}\right] \] which have the property that
all operators along a row (or column) commute. Further,
the product of rows (or columns) yields
$\hat{R}_{i}=\mathbb{I}$ and
$\hat{C}_{j}=\mathbb{I}\,(j\neq3)$,
$\hat{C}_{3}=-\mathbb{I}$, ($\forall\,i,j$) where
$\hat{R}_{i}:=\prod_{j}\hat{A}_{ij}$,
$\hat{C}_{j}:=\prod_{i}\hat{A}_{ij}$. Let us
assume that $M$ exists. From property (b) of the
map, to get $M(\hat{C}_{3})=-1$ (as required by
property (a)), we must have an odd number of $-1$
assignments in the third column. In the remaining
columns, the number of $-1$ assignments must be
even (for each column). Thus, in the entire
square, the number of $-1$ assignments must be
odd. Let us use the same reasoning, but along the
rows. Since each $M(\hat{R}_{i})=1$, we must have
even number of $-1$ assignments along each row.
Thus, in the entire square, the number of $-1$
assignments must be even. We have arrived at a
contradiction and therefore we conclude that our
assumption that a consistent $M$ exists, must be
wrong.  \end{proof}

\begin{thmrem} One could in principle assume $M$, to
be s.t. (a) $M(\hat{\mathbb{I}})=1$, (b)
$M(\alpha\hat{B}_{i})=\alpha M(\hat{B}_{i})$, for
$\alpha\in\mathbb{R}$, (c)
$M(\hat{B}_{i}^{2})=M(\hat{B}_{i})^{2}$, (d)
$M(\hat{B}_{i}+\hat{B}_{j})=M(\hat{B}_{i})+M(\hat{B}_{j})$,
to deduce (d')
$M(\hat{B}_{i}\hat{B}_{j})=M(\hat{B}_{i})M(\hat{B}_{j})$
and that $M(\hat{B}_{i})\in$ spectrum of
$\hat{B}_{i}$. Effectively then, condition (b)
listed in the theorem is satisfied as a
consequence.  Therefore, assuming (a)-(d) as
listed here, rules out a larger class of $M$.
\cite{KochenSpecker} 
\end{thmrem} 
Here $M$ maybe viewed as a specific
class of prediction maps, that implicitly depends
on the state $[\psi]$. It is clear that according to
the theorem, non-contextual maps which are
\emph{multiplicative} must be incompatible with
QM. However, non-contextual maps which are
non-\emph{multiplicative} might be consistent with
QM. Indeed it is and we will show this by constructing an
explicit model.


Before proceeding we note, however, that QM can
enforce \emph{sequential multiplicativity}. This
result will be used to assess accuracy of any
proposal of a HV
model.
\begin{prop*} 
Let a quantum mechanical 
system be in a state, s.t. measurement of
$\hat{C}$ yields repeatable results (same result
each time). Then according to QM, \emph{sequential
multiplicativity} holds, where $\hat{C}:=
f(\hat{B}_{1},\hat{B}_{2},\dots\hat{B}_{n})$, and
$\hat{B}_{i}$ are as defined (in \defnref{seqnMltpl})
\end{prop*}
\begin{proof} 
Without loss of generality we can take
$\hat{B}_{1},\hat{B}_{2},\dots\hat{B}_{n}$ to be
mutually compatible and a complete set of
observables (operators can be added to make the set
complete if it not so).
It follows
that $\exists$
$\left|\mathbf{b}=\left(b_{1},b_{2},\dots
b_{n}\right)\right\rangle $ s.t.
$\hat{B}_{i}\left|\mathbf{b}\right\rangle
=b_{i}\left|\mathbf{b}\right\rangle $,
and that
$\sum_{\mathbf{b}}\left|\mathbf{b}\right\rangle
\left\langle \mathbf{b}\right|=\hat{\mathbb{I}}$.
Let the state of the system
$\left|\psi\right\rangle $ be s.t.
$\hat{C}\left|\psi\right\rangle
=c\left|\psi\right\rangle $. For
the statement to follow, one need only show that
$\left|\psi\right\rangle $ must be made of only
those $\left|\mathbf{b}\right\rangle $s, which
satisfy $c=\tilde
f(b_{1},b_{2},\dots
b_{n})$.  This is the crucial step and
proving this is straightforward. We start with
$\hat{C}\left|\psi\right\rangle
=c\left|\psi\right\rangle $ and take its inner
product with $\left\langle \mathbf{b}\right|$ to
get \begin{eqnarray*} \left\langle
\mathbf{b}\right|\hat{C}\left|\psi\right\rangle  &
= & c\left\langle \mathbf{b}|\psi\right\rangle ,\\
\left\langle
\mathbf{b}\right|f(\hat{B}_{1},\hat{B}_{2},\dots\hat{B}_{n})
\left|\psi\right\rangle
& = & c\left\langle \mathbf{b}|\psi\right\rangle
,\\ \tilde f(b_{1},b_{2},\dots
b_{n})\left\langle
\mathbf{b}|\psi\right\rangle  & = & c\left\langle
\mathbf{b}|\psi\right\rangle .  \end{eqnarray*}
Also, we have $\left|\psi\right\rangle
=\sum_{\mathbf{b}}\left\langle
\mathbf{b}|\psi\right\rangle
\left|\mathbf{b}\right\rangle $, from
completeness. If we consider
$\left|\mathbf{b}\right\rangle $s for which
$\left\langle \mathbf{b}|\psi\right\rangle \neq0$,
then we can conclude that indeed $c=\tilde
f(b_{1},b_{2},\dots
b_{n})$. However, when $\left\langle
\mathbf{b}|\psi\right\rangle =0$, viz.
$\left|\mathbf{b}\right\rangle $s that are
orthogonal to $\left|\psi\right\rangle $, then
nothing can be said but that is irrelevant. We can thus conclude that
$\left|\psi\right\rangle $ is made only of those
$\left|\mathbf{b}\right\rangle $s that satisfy the
required relation.
\end{proof} 
It is worth noting that in the Peres Mermin
case, where $\hat{R}_{i}$ and $\hat{C}_{j}$ are
just $\pm\hat{\mathbb{I}}$, it follows that all
states are their eigenstates. Consequently, for
these operators \emph{sequential multiplicativity}
must always hold.
%\subsection{C-ingle model | An Explicit
%Construction \label{ssec:cingle}}
We are now ready to describe our explict model.
Let the state of a finite dimensional quantum system
be $\left|\chi\right\rangle $. We wish to
assign a value to an arbitrary observable 
$\hat{A}=\sum_{a}a\left|a\right\rangle
\left\langle a\right|$, which has 
eigenvectors $\{ \vert
a_j\rangle \}$. The corresponding ordered eigenvalues are $\{a_j\}$ such that 
$a_{\rm min}=a_1$
and 
$a_{\rm max}=a_n$.

Our HV model for QM assigns values in the following
three steps:
\setdefaultleftmargin{0pt}{}{}{}{}{}
\begin{enumerate}
\item
{\bf Initial HV:} Pick a number
$c\in[0,1]$, from a uniform random distribution.\
\item
{\bf
Assignment or Prediction:}
 The value assigned to
$\hat{A}$ is given by finding the smallest $a$
s.t.  $c\le\sum_{a'=a_{\text{min}}}^{a}\left|\left\langle
a'|\chi\right\rangle \right|^{2},$ viz. we have
specified a prediction map, $M(\hat{A},[\chi ])=a$.
\item
{\bf Update:} After measuring an operator, the state is
updated (collapsed) in accordance with the rules
of QM. This completely specifies the sequence map
$S$.
\end{enumerate}
To see how this works consider 
a spin-half particle in the state
$\left|\chi\right\rangle
=\cos\theta\left|0\right\rangle
+\sin\theta\left|1\right\rangle $ and the observable
$\hat{A}=\hat{\sigma}_{z}=\left|0\right\rangle
\left\langle 0\right|-\left|1\right\rangle
\left\langle 1\right|$. Now, according to the
postulates of this theory, $M(\hat{A},[\chi])=+1$, if
the randomly generated value for 
$c\le\cos^{2}\theta$ else $\hat{A}$ is assigned
$-1$. It follows then, from $c$ being uniformly
random in $[0,1]$ that the statistics agree with
predictions of QM.

Note that the assignment, described by the
prediction map $M$, is non-contextual since given
an operator and a state (+ the HV $c$), the value is
uniquely assigned. The map $M$ is, however,
non-multiplicative. 


To bring out the non-multiplicativity of our model 
and to convince the reader of its generic validity,
we apply it to the Peres Mermin situation of two
spin-half particles. Consider the initial state of
the system 
$\left|\psi_{1}\right\rangle
=\left|00\right\rangle $. Assume we 
obtained $c=0.4$ as a random choice. To arrive
at the assignments, note that
$\left|00\right\rangle $ is an eigenket of only
$\hat{R}_{i},\hat{C}_{j}$ and
$\hat{A}_{33}=\hat{\sigma}_{z}\otimes\hat{\sigma}_{z}$.
Thus, in the first iteration, all these should be
assigned their respective eigenvalues. The
remaining operators must be assigned $-1$ as one
can readily verify by explicitly finding the
smallest $a$ as described in postulate 2 of the
model (see \tblref{HVmodel}). Two remarks are in
order. First, this model is manifestly
\emph{non-multiplicative}, for
$M_{1}(\hat{C}_{3})=1\neq
M_{1}(\hat{A}_{13})M_{1}(\hat{A}_{23})M_{1}(\hat{A}_{33})=-1$,
where the subscript has been introduced to index
the iteration. More precisely,
$M_{1}(\hat{O}):=M(\hat{O},\left[\left|\psi_{1}\right\rangle
=\left|00\right\rangle \right])$ where the
complete state $\left[\left|\psi_{1}\right\rangle
\right]$ implicitly refers to both the quantum
state $\left|00\right\rangle $ and the HV $c=0.4$.
Second, we must impose \emph{sequential
multiplicativity} as a consistency check of the
model, which in particular entails that
$M_{1}(\hat{C}_{3})=M_{1}(\hat{A}_{33})M_{2}(\hat{A}_{23})M_{3}(\hat{A}_{13})$,
where
$M_{2}:=M(\hat{O},\left[\left|\psi_{2}\right\rangle
\right])$,
$M_{3}:=M(\hat{O},\left[\left|\psi_{3}\right\rangle
\right])$ and $\left|\psi_{2}\right\rangle
,\,\left|\psi_{3}\right\rangle $ are obtained from
postulate 3. Note that for each iteration, a new
HV is generated. To illustrate \emph{sequential
multiplicativity}, we must choose to measure
$\hat{A}_{33}$. According to step 3, since
$\left|00\right\rangle $ is an eigenstate of
$\hat{A}_{33}$, the final state remains
$\left|00\right\rangle $.  


For the next iteration, $i=2$, viz.
after the first measurement has been performed,
say the random number generator yielded $c=0.1$.
Since $\left|\psi\right\rangle $ is also unchanged
the assignment remains invariant (in fact any of
$c<0.5$ would yield the same result), which should
be evident from the previous exercise. For the
final step we choose to measure
$\hat p := \hat{A}_{23}(=\hat{\sigma}_{y}\otimes\hat{\sigma}_{y})$,
to proceed with sequentially measuring
$\hat{C}_{3}$. To simplify calculations, we note
\[ \left|00\right\rangle
=\frac{(\left|\tilde{+}\tilde{-}\right\rangle
+\left|\tilde{-}\tilde{+}\right\rangle
)/\sqrt{2}+(\left|\tilde{+}\tilde{+}\right\rangle
+\left|\tilde{-}\tilde{-}\right\rangle
)/\sqrt{2}}{\sqrt{2}}, \] where
$\left|\tilde{\pm}\right\rangle
=\left|0\right\rangle \pm i\left|1\right\rangle $
(eigenkets of $\hat{\sigma}_{y}$). Since
$\left|00\right\rangle $ is manifestly not an
eigenket of $\hat{p}$, we must find an appropriate
eigenket $\left|p_{-}\right\rangle $ s.t.
$\hat{p}\left|p_{-}\right\rangle
=-\left|p_{-}\right\rangle $, since $c=0.1$ and
$\left\langle p_{-}|00\right\rangle $ is already
$>0.1$. It is immediate that
$\left|p_{-}\right\rangle
=\left(\left|\tilde{+}\tilde{-}\right\rangle
+\left|\tilde{-}\tilde{+}\right\rangle
\right)/\sqrt{2}=\left(\left|00\right\rangle
+\left|11\right\rangle \right)/\sqrt{2}$, which
becomes the final state. \\ For the final
iteration, $i=3$, say we obtain $c=0.7$. So far,
we have $M_{1}(\hat{A}_{33})=1$ and
$M_{2}(\hat{A}_{23})=-1$. We must obtain
$M_{3}(\hat{A}_{13})=1$, independent of the value
of $c$, to be consistent. Let's check that.
Indeed, according to postulate 2, since
$\hat{\sigma}_{x}\otimes\hat{\sigma}_{x}\left(\left|00\right\rangle
+\left|11\right\rangle
\right)/\sqrt{2}=1\left(\left|00\right\rangle
+\left|11\right\rangle \right)/\sqrt{2}$,
$M_{3}(\hat{A}_{13})=1$ for all allowed values of
$c$. As a remark, it maybe emphasised that the
$M_{2}(\hat{A}_{33})=M_{3}(\hat{A}_{33})$ and
$M_{2}(\hat{A}_{23})=M_{3}(\hat{A}_{23})$, which
essentially expresses compatibility of these
observables, viz. measurement of $\hat{A}_{13}$
doesn't affect the result one would obtain by
measuring operators compatible to it (granted they
have been measured once before).

%\begin{widetext}
%\onecolumngrid
\begin{table*}
 \begin{equation*}
\begin{array}{ccc} 
i=1,c=0.4,\left|\psi_{\text{init}}\right\rangle=\left|00\right\rangle & i=2,c=0.1,\left|\psi_{\text{init}}\right\rangle=\left|00\right\rangle & i=3,c=0.7,\left|\psi_{\text{init}}\right\rangle=\frac{\left|00\right\rangle +\left|11\right\rangle}{\sqrt{2}} \\
%\text{Iteration} & i=1 & i=2 & i=3\\ 
%\left|\psi_{\text{init}}\right\rangle  & \left|00\right\rangle  & \left|00\right\rangle  & \frac{\left|00\right\rangle +\left|11\right\rangle}{\sqrt{2}}\\ 
%\text{HV} & c=0.4 & c=0.1 & c=0.7\\ \\ 

M_{1}(\hat{A}_{ij})\doteq\left[\begin{array}{ccc} -1 & -1 & -1\\ -1 & -1 & -1\\ -1 & -1 & +1\end{array}\right] &
M_{2}(\hat{A}_{ij})\doteq\left[\begin{array}{ccc}-1 & -1 & -1\\ -1 & -1 & -1\\ -1 & -1 & +1\end{array}\right] &
M_{3}(\hat{A}_{ij})\doteq\left[\begin{array}{ccc}+1 & +1 & +1\\ +1 & +1 & -1\\ +1 & +1 & +1\end{array}\right]\\ \\


M_{1}(\hat{R}_{i}),M_{1}(\hat{C}_{j})=+1\,(j\neq3) &
M_{2}(\hat{R}_{i}),M_{2}(\hat{C}_{j})=+1\,(j\neq3) &
M_{3}(\hat{R}_{i}),M_{3}(\hat{C}_{j})=+1\,(j\neq3)\\

M_{1}(\hat{C}_{3})=-1 & M_{2}(\hat{C}_{3})=-1 & M_{3}(\hat{C}_{3})=-1\\ 


\hat{A}_{13}=\hat{\sigma}_{z}\otimes\hat{\sigma}_{z};M_{1}(\hat{A}_{13})=+1&
\quad\quad\hat{A}_{23}=\hat{\sigma}_{y}\otimes\hat{\sigma}_{y};M_{2}(\hat{A}_{23})=-1\quad\quad&
\hat{A}_{33}=\hat{\sigma}_{x}\otimes\hat{\sigma}_{x};M_{3}(\hat{A}_{33})=+1\\

\left|\psi_{\text{final}}\right\rangle = \left|00\right\rangle  & \left|\psi_{\text{final}}\right\rangle = \frac{\left|00\right\rangle +\left|11\right\rangle}{\sqrt{2}} & \left|\psi_{\text{final}}\right\rangle = \frac{\left|00\right\rangle+\left|11\right\rangle }{\sqrt{2}}
% \begin{bmatrixT} 
% \text{Iteration} & i=1 & i=2
% & i=3\\ \left|\psi_{\text{init}}\right\rangle  &
% \left|00\right\rangle  & \left|00\right\rangle  &
% \frac{\left|00\right\rangle +\left|11\right\rangle
% }{\sqrt{2}} \\
% \text{HV} & c=0.4 & c=0.1 & c=0.7 \\

% \text{Predictions or} &
% {m_{1}(\hat{A}_{ij})\doteq\left[\begin{array}{ccc}
% -1 & -1 & -1\\ -1 & -1 & -1\\ -1 & -1 & +1
% \end{array}\right]} &
% {m_{2}(\hat{A}_{ij})\doteq\left[\begin{array}{ccc}
% -1 & -1 & -1\\ -1 & -1 & -1\\ -1 & -1 & +1
% \end{array}\right]} &
% {m_{3}(\hat{A}_{ij})\doteq\left[\begin{array}{ccc}
% +1 & +1 & +1\\ +1 & +1 & -1\\ +1 & +1 & +1
% \end{array}\right]}\\ 

% \text{(Assignments)} &
% m_{1}(\hat{R}_{i}),m_{1}(\hat{C}_{j})=+1\,(j\neq3)
% &
% m_{2}(\hat{R}_{i}),m_{2}(\hat{C}_{j})=+1\,(j\neq3)
% &
% m_{3}(\hat{R}_{i}),m_{3}(\hat{C}_{j})=+1\,(j\neq3)\\
% & m_{1}(\hat{C}_{3})=-1 & m_{2}(\hat{C}_{3})=-1 &
% m_{3}(\hat{C}_{3})=-1
% \end{bmatrixT}
\label{eq:toyModel}
\end{array}
% \label{eq:toyModel}
\end{equation*}
\caption{HV model applied to the Peres Mermin situation}
\label{tbl:HVmodel}
\end{table*}
%\end{widetext} 




%%%%%%%%%%%%%%%%%%%%%%%%%%%%%%%%%%%
{\color{red} Will work on CC after everthin else
is done}
This has implications to Bohm's HV model as well, Bohmian Mechanics (BM), where the notion of positions (and hence trajectories) is precisely defined and particles are guided by the wavefunction. Consequently, given the initial conditions precisely (which is never the case experimentally) one can predict results of all experiments precisely. This entails that BM is capable of predicting outcomes of, for instance, a Stern Gerlach experiment. However, within BM two experiments which correspond to measurement of the same operator, can yield opposite results. Further, this can be demonstrated for a single spin half particle. This observation is often used as a demonstration of `contextuality' in BM, although we did not even speak of compatible measurements. It turns out, however, that if one restricts oneself to a one-one map between an experiment and an observable (by following some convention for example) then the so called `contextuality' can be removed from BM. One finds that the map so obtained from observables to their predictions turns out to be non-multiplicative, suggesting that at least in this case, non-multiplicativity is a more suitable explanation of non-classicality.

We have already constructed an explicit non-contextual model,
which is consistent with QM.  This model we knew
had to be non-multiplicative.  We will see how
non-multiplicativity gives rise to what one might
confuse to mean contextuality.  Imagine $\hat{B}_{1}=\hat{\sigma}_{z}\otimes\hat{\mathbb{I}}=\left|00\right\rangle
\left\langle 00\right|+\left|01\right\rangle
\left\langle 01\right|-\left[\left|10\right\rangle
\left\langle 10\right|+\left|11\right\rangle
\left\langle 11\right|\right],$ $
\hat{B}_{2}=\hat{\mathbb{I}}\otimes\hat{\sigma}_{z}=\left|10\right\rangle
\left\langle 10\right|+\left|11\right\rangle
\left\langle 11\right|-\left[\left|00\right\rangle
\left\langle 00\right|+\left|01\right\rangle
\left\langle 01\right|\right],$ while we define
$ \hat{C}=f(\{\hat{B}_{i}\})
=0.\left|00\right\rangle \left\langle
00\right|+1.\left|01\right\rangle \left\langle
01\right|+2.\left|10\right\rangle \left\langle
10\right|+3.\left|11\right\rangle \left\langle
11\right|. $
% \begin{align*} \hat{C}&=f(\{\hat{B}_{i}\}) \\
% &=0.\left|00\right\rangle \left\langle
% 00\right|+1.\left|01\right\rangle \left\langle
% 01\right|+2.\left|10\right\rangle \left\langle
% 10\right|+3.\left|11\right\rangle \left\langle
% 11\right|.  \end{align*}
 $\hat{C}$ maybe viewed as
a function of $\hat{B}_{1}$, $\hat{B}_{2}$ and
other operators $\hat{B}_{i}$ which are
constructed to obtain a maximally commuting set. A
measurement of $\hat{C}$, will collapse the state
into one of the states which are simultaneous
eignkets of $B_{1}$ and $B_{2}$. Consequently,
from the observed value of $\hat{C}$, one can
deduce the values of $\hat{B}_{1}$ and
$\hat{B}_{2}$. Now consider
$\sqrt{2}\left|\chi\right\rangle
=\left|10\right\rangle +\left|01\right\rangle $,
for which $M_{1}(\hat{B}_{1})=1$, and
$M_{1}(\hat{B}_{2})=1$, using the C-ingle model,
with $c<0.5$. However, $M_{1}(\hat{C})=1$, from
which one can deduce that $B_{1}$ was $+1$, while
$B_{2}$ was $-1$.  This property itself, one may
be tempted call contextuality, viz.  the value of
$B_{1}$ depends on whether it is measured alone or
with the remaining $\{B_{i}\}$. However, it must
be noted that $B_{1}$ has a well defined value,
and so does $\hat{C}$. Thus by our accepted
definition, there's no contextuality. It is just
that $M_{1}(\hat{C})\neq
\tilde f(M_{1}(\hat{B}_{1}),M_{1}(\hat{B}_{2}),\dots)$,
viz. the theory is non-multiplicative. Note that
after measuring $\hat{C}$ however,
$M_{2}(\hat{B}_{1})=+1$ and
$M_{2}(\hat{B}_{2})=-1$ (for any value of $c$)
consistent with those deduced by measuring
$\hat{C}$.


It is known that for simultaneous eigenstates of
$\hat{B}_{i}$, multiplicativity must hold.
Consequently, any violation of multiplicativity
must arise from states that are superpositions (of
the simultaneous eigenkets).  Note, however, that
entanglement is not necessary to show a violation
since, for example, the PM test is a state
independent test where a separable state can be
used to arrive at a contradiction. Entanglement 
can be used to show non-locality of non-multiplicative
prediction maps using Bell's theorem.
\bibliographystyle{apsrev4-1}
\bibliography{references}

\end{document}
